\documentclass[11pt, a4paper]{article}

\usepackage{amsmath}
\usepackage{amssymb}
\usepackage{geometry}
\usepackage{hyperref}
\usepackage{enumitem}
\usepackage{booktabs}
\usepackage{tcolorbox}
\usepackage{microtype}
\usepackage{parskip}
\usepackage{titlesec}
\usepackage{fancyhdr}
\usepackage{xcolor}

\geometry{
  a4paper,
  top=3cm, bottom=3cm,
  left=3cm, right=3cm
}

% ── Colours ──────────────────────────────────────────────────────────────────
\definecolor{inksoft}{RGB}{74, 71, 64}
\definecolor{inkmuted}{RGB}{138, 134, 128}
\definecolor{accent}{RGB}{45, 74, 138}
\definecolor{gold}{RGB}{138, 106, 26}
\definecolor{rulecolor}{RGB}{216, 210, 196}
\colorlet{RULECOLOR}{rulecolor}
\definecolor{boxbg}{RGB}{242, 237, 227}
\definecolor{accentbg}{RGB}{232, 237, 247}
\definecolor{goldbg}{RGB}{245, 240, 224}

% ── Section formatting ────────────────────────────────────────────────────────
\titleformat{\section}
  {\small\bfseries\color{inkmuted}\ttfamily\uppercase}
  {}{0em}{}[\vspace{2pt}\noindent\textcolor{rulecolor}{\rule{\linewidth}{0.4pt}}\vspace{2pt}]

\titlespacing{\section}{0pt}{2em}{0.75em}

% ── Header / footer ───────────────────────────────────────────────────────────
\pagestyle{fancy}
\fancyhf{}
\renewcommand{\headrulewidth}{0pt}
\fancyfoot[C]{\small\color{inkmuted}\thepage}

% ── Hyperref ──────────────────────────────────────────────────────────────────
\hypersetup{
  colorlinks=true,
  linkcolor=accent,
  urlcolor=accent,
  pdfauthor={Juha Meskanen},
  pdftitle={The Universe Lives Inside You — An Introduction to IAME}
}

% ── Custom environments ───────────────────────────────────────────────────────
\tcbuselibrary{skins, breakable}

\newtcolorbox{iepbox}{
  colback=accentbg,
  colframe=accent,
  leftrule=3pt, rightrule=0pt, toprule=0pt, bottomrule=0pt,
  arc=2pt,
  left=10pt, right=10pt, top=8pt, bottom=8pt,
  fontupper=\small
}

\newtcolorbox{axiombox}{
  colback=boxbg,
  colframe=rulecolor,
  leftrule=0.5pt, rightrule=0.5pt, toprule=0.5pt, bottomrule=0.5pt,
  arc=2pt,
  left=12pt, right=12pt, top=10pt, bottom=10pt
}

\newtcolorbox{openbox}{
  colback=goldbg,
  colframe=gold,
  leftrule=0.5pt, rightrule=0.5pt, toprule=0.5pt, bottomrule=0.5pt,
  arc=2pt,
  left=12pt, right=12pt, top=10pt, bottom=10pt,
  fontupper=\small
}

\newtcolorbox{pullbox}{
  colback=white,
  colframe=white,
  left=20pt, right=20pt, top=4pt, bottom=4pt,
  fontupper=\large\itshape\color{inksoft}
}

% ── Misc ──────────────────────────────────────────────────────────────────────
\newcommand{\iame}{\textsc{iame}}
\newcommand{\Cs}{\mathcal{C}_s}
\newcommand{\Cg}{\mathcal{C}_g}
\newcommand{\Linfo}{\mathcal{L}_{\mathrm{info}}}

\setlength{\parskip}{0.65em}
\setlength{\parindent}{0pt}

% ─────────────────────────────────────────────────────────────────────────────
\begin{document}

% ── Masthead ─────────────────────────────────────────────────────────────────
\noindent{\color{accent}\rule{\linewidth}{1.5pt}}

\vspace{0.5em}

{\small\ttfamily\color{inkmuted} Introduction · \iame{} Framework}

\vspace{0.4em}

{\fontsize{26}{32}\selectfont\bfseries The Universe Lives Inside You}

\vspace{0.3em}

{\large\itshape\color{inksoft}
An introduction to the Internal Emergence framework --- where
consciousness, physics, and mathematics arise from a single principle}

\vspace{0.8em}

{\small\ttfamily\color{inkmuted}
Juha Meskanen \quad·\quad Papers I--IV \quad·\quad 2001--2025}

\vspace{0.5em}

\noindent{\color{rulecolor}\rule{\linewidth}{0.5pt}}

\vspace{1.5em}

{\itshape\color{inksoft}
What if the deepest question in physics --- why does the universe look the
way it does? --- has the same answer as the deepest question in
consciousness --- why is there something it is like to be me? The \iame{}
framework says: because they are the same question.}

\vspace{1.5em}

% ─────────────────────────────────────────────────────────────────────────────
\section{Starting from nothing}

Every theory of everything starts somewhere. String theory starts with
extended objects vibrating in a background spacetime. Loop quantum gravity
starts with a discrete network of quantum geometry. The Standard Model starts
with a specific list of particles and symmetries. Each of these frameworks
imports considerable structure before it derives anything.

The \iame{} framework starts with less. It starts with the only thing that
cannot be coherently denied: \emph{something is experiencing}. You cannot
argue against this without performing the very act of experience that the
argument denies. This is not a philosophical claim --- it is the one piece
of data that requires no inference and no external measurement.

From this single starting point, and from a handful of modest empirical
observations about biology and physics, the framework derives the structure
of the observed universe: why time flows, why space expands, why there are
particles rather than nothing, why quantum mechanics and general relativity
look the way they do --- and why mathematics is unreasonably effective at
describing all of it.

% ─────────────────────────────────────────────────────────────────────────────
\section{The axiomatic chain}

The argument proceeds as a chain of six observations, each individually
modest, collectively powerful.

\begin{axiombox}
{\small\ttfamily\color{inkmuted} Six premises · each independently rejectable}

\vspace{0.75em}

\begin{description}[leftmargin=2.5em, labelwidth=2em, labelsep=0.5em,
                    itemsep=0.35em, parsep=0pt]
  \item[\textcolor{accent}{\small A1}]
    Something is experiencing. The denial of this is itself an experience.
  \item[\textcolor{accent}{\small A2}]
    The observer is finite --- bounded memory, bounded lifespan.
  \item[\textcolor{accent}{\small A3}]
    DNA is the only mechanism we have ever observed that reliably produces
    conscious observers.
  \item[\textcolor{accent}{\small A4}]
    DNA is ordinary matter obeying ordinary physical laws. No magic.
  \item[\textcolor{accent}{\small A5}]
    The physical processes underlying DNA are computable --- in principle,
    simulable.
  \item[\textcolor{accent}{\small A6}]
    Pain is not epiphenomenal. It has measurable causal effects on behaviour.
\end{description}

\vspace{0.5em}
\noindent\textcolor{inkmuted}{$\downarrow$}
\vspace{0.3em}

\noindent\textcolor{gold}{$\therefore$}\enspace
\textbf{A conscious observer is a finite bitstring.} A sufficiently large
collection of bitstrings contains one --- with no running process, no
external reader, no simulation required.
\end{axiombox}

The key step is A5 combined with A6. If DNA processes are computable and
pain has real causal effects, then a perfect simulation of a pain-sensitive
human must itself experience pain --- because if it did not, the experiential
difference would produce a behavioural difference, contradicting the
assumption of perfect simulation. This is not a philosophical argument about
zombies. It is a logical consequence of the axioms, and it is refutable:
build a sufficiently accurate DNA simulation and demonstrate that it lacks
pain responses.

Once a conscious observer is established as a finite bitstring, the next
step follows from a thought experiment. Take that running simulation and
progressively replace computed state transitions with precomputed lookup
tables. At each step, the sequence of internal states --- everything Alice
experiences --- is identical. In the limit, all computation is replaced by a
static, precomputed dataset. If consciousness vanished at some point during
this optimisation, there would have to be a minimum code-to-data ratio below
which experience ceases. That ratio would be a new physical constant with no
basis in any of the six axioms. Its existence is therefore excluded. Alice's
experience is a property of the information structure itself, not of any
process that generates it.

\begin{iepbox}
\textbf{The Internal Emergence Principle.} Consciousness is a property of
certain finite information structures, not of processes acting upon them.
Any configuration space large enough to contain such structures contains
consciousness --- with no simulation, no clock, and no external observer
required.
\end{iepbox}

% ─────────────────────────────────────────────────────────────────────────────
\section{What follows immediately}

The Internal Emergence Principle has consequences that are stark and
non-negotiable.

\textbf{Time is not fundamental.} A conscious observer exists within a
static bitstring --- no external clock is running. Time, as experienced,
must therefore be defined by internal structure alone. The only available
ordering is: configuration $B$ follows configuration $A$ if and only if $B$
contains a memory of $A$. The flow of time is an ordinal index internal to
the observer. The universe, at the level of the totality, is static. This
conclusion is independently supported by the block-universe picture of
general relativity, in which the four-dimensional spacetime manifold is a
static geometric object and the field equations contain no preferred
direction of time.

\textbf{The observer and the universe are identical.} In a static universe
with no running external process, no mechanism exists by which external bits
could alter internal bits. Every structure Alice observes --- every star,
every sensation, every other person --- must already be encoded within the
bitstring that constitutes her. Alice does not receive information from
outside. She discovers structure already present within her. It follows that
Alice is not a small observer inside a vast universe. The vast universe she
experiences is her interior. And since Alice is part of her universe
(trivially) and her universe is part of Alice (by the containment result),
basic set theory gives:

\begin{equation}
  \boxed{O = U_O}
\end{equation}

For every observer $O$, the observer and the universe it experiences are
identical structures. What two observers share is the intersection of their
universes --- the compressible, verifiable physics that constitutes
intersubjective reality. There is no view from nowhere.

% ─────────────────────────────────────────────────────────────────────────────
\section{The measure: what is typical}

Establishing that a static configuration space contains conscious observers
is not enough. We also need to know which configurations a conscious observer
is likely to find themselves in. The answer comes from a compression measure.

Under any Solomonoff-style measure, configurations with shorter descriptions
are exponentially more probable than configurations with longer ones. A
universe of smooth, lawlike, inertial physics has a short description. A
universe of random chaos has a long description and is exponentially
suppressed. An observer is therefore \emph{expected} --- not lucky --- to
find themselves in a smooth, structured, lawlike universe. Boltzmann brains,
chaotic fluctuations, and arbitrary random configurations are all disfavoured
by the same measure that selects physical structure. This is not fine-tuning.
It is counting.

The specific form of the compression measure that the framework identifies
has two components:

\begin{equation}
  \Linfo = \Cs + \Cg \;\longrightarrow\; \text{minimum}
\end{equation}

$\Cs$ is the \emph{spectral complexity} --- the bit-cost of encoding
wave-like, relational, quantum-mechanical behaviour. $\Cg$ is the
\emph{geometric complexity} --- the bit-cost of encoding locality,
boundaries, and large-scale spatial structure. Observed physics is the
configuration that minimises both simultaneously.

\vspace{0.5em}

\noindent\begin{tabular}{p{0.3\linewidth} p{0.3\linewidth} p{0.3\linewidth}}
\toprule
\textbf{D} \quad Data / Particles &
\textbf{$\psi$} \quad Spectral / Quantum &
\textbf{G} \quad Geometry / Gravity \\
\midrule
Discrete, localised digital information. The finite output when a bounded
observer parses the joint minimum. &
The smooth wave-codec that compresses relational possibility space.
Minimises $\Cs$. This is quantum mechanics. &
The continuous metric manifold that compresses realisation space at scale.
Minimises $\Cg$. This is general relativity. \\
\bottomrule
\end{tabular}

\vspace{0.75em}

Quantum mechanics and general relativity are not two independent theories in
need of unification. They are the two complementary compression schemes that
minimise total description length for a finite observer. Their apparent
incompatibility is not a deep physical mystery --- it is an artefact of
treating them as fundamental rather than as emergent tools.

% ─────────────────────────────────────────────────────────────────────────────
\section{Singularities and expansion}

The representational equivalence between a physical system and its
computational simulation has an immediate consequence: properties that are
geometrically inaccessible --- regions where the spacetime manifold breaks
down --- remain accessible through the execution trace of the simulation.
When general relativity fails, the bitstring does not.

Applying this to gravitational collapse, the Shannon entropy of the geometric
encoding of the execution trace decreases monotonically as particles fall
toward the singularity, extrapolating to zero at the classical endpoint. A
zero-entropy bitstring encodes a single, structureless state --- a geometric
point. It cannot support particles, microstructure, or stress-energy. The
gravitational singularity is not a catastrophic breakdown of physics. It is
an informational fixed point of zero complexity, the simplest state the
universe can be in.

\begin{pullbox}
The singularity is not where physics breaks down.\\
It is where physics becomes trivially simple.
\end{pullbox}

Now reverse the process. Begin from that zero-entropy state --- the unique
all-zero bitstring, the simplest possible configuration --- and increase
entropy through random mutations. Three phenomena emerge without any
hard-coded physical laws, equations of motion, or imposed initial
irregularities:

\begin{enumerate}[label=\textcolor{gold}{\small\texttt{\arabic*.}},
                  leftmargin=2.5em, itemsep=0.4em]
  \item \textbf{Exponential spatial expansion} from a structureless point.
        Cosmic expansion is the geometric reading of entropy increase ---
        not a dynamical process imposed on space.
  \item \textbf{Hierarchical particle-like structures} whose abundance follows
        a lognormal distribution at each entropy level --- particles first,
        then composite objects, then higher-order structures.
  \item \textbf{Scale invariance} of particle sizes under global expansion,
        exactly as observed.
\end{enumerate}

These results hold across more than two hundred distinct structural filter
definitions, geometric encoding schemes, and representational levels ---
including purely combinatorial pattern-matching on raw bitstrings with no
geometric embedding whatsoever. The lognormal distribution of emerging
structure is not a property of any particular physical model. It is an
intrinsic combinatorial property of entropy-increasing bitstrings. The arrow
of time, the thermodynamic second law, and the expansion of the universe are
not three independent features of the observed world. They are one feature,
seen through three different lenses.

% ─────────────────────────────────────────────────────────────────────────────
\section{The Dirac conjecture}

The two-term structure $\Linfo = \Cs + \Cg$ raises an immediate question:
are these genuinely independent quantities, or are they two frequency regimes
of a single deeper object?

A foundational result of spectral geometry establishes that the spectrum of
the Dirac operator $D$ on a Riemannian manifold encodes its full spin
geometry --- metric, orientation, and curvature are all recoverable from the
eigenvalue spectrum. Connes' spectral action principle defines a natural
complexity functional on geometries described by $D$:

\begin{equation}
  S_{\mathrm{spectral}}
  = \mathrm{Tr}\!\left(f\!\left(\frac{D}{\Lambda}\right)\right)
\end{equation}

When this functional is expanded at low frequencies, it recovers the
Einstein-Hilbert action of general relativity. When evaluated at high
frequencies, it encodes wave-like quantum behaviour. The conjecture of the
\iame{} framework is that $\Cg$ is the low-frequency limit of this
functional and $\Cs$ is its high-frequency residual --- so that the
informational Lagrangian reduces to a single spectral action built from one
operator:

\begin{equation}
  \boxed{
    \Linfo
    = \mathrm{Tr}\!\left(f\!\left(\frac{D}{\Lambda}\right)\right)
    \;\longrightarrow\; \text{minimum}
  }
\end{equation}

If confirmed, quantum mechanics and general relativity are not two theories.
They are the infrared and ultraviolet limits of minimising one expression.
The Dirac operator wins not because it is postulated but because it is the
minimum-description-length encoding of spin geometry --- it is what the
compression measure selects when no operator is chosen by hand.

This conjecture is the primary open mathematical problem of the framework.
It is well-posed and awaits verification in spectral geometry and
noncommutative geometry.

% ─────────────────────────────────────────────────────────────────────────────
\section{Falsifiable predictions}

\begin{enumerate}[label=\textcolor{gold}{\small\texttt{P\arabic*.}},
                  leftmargin=2.5em, itemsep=0.5em]
  \item Any stable observer-hosting physics must exhibit both wave-like and
        geometric structure. Pure geometry without quantum fields, or pure
        waves without geometry, is forbidden.
  \item Discreteness is never fundamental. Any deeper physical theory will
        not be fundamentally discrete. Discreteness emerges at the observer
        boundary.
  \item The correct theory of quantum gravity will take the form of a joint
        spectral-geometric minimisation --- specifically, a Dirac spectral
        action --- rather than a quantisation of a background metric.
  \item The lognormal microstructure distribution near gravitational
        singularities is universal across geometries, filters, and
        representational levels. The shape parameters at the self-consistent
        filter correspond to physically observable quantities.
  \item The Pauli exclusion principle emerges from description-length
        explosion upon observer merger, without antisymmetry being
        hard-coded.
  \item A smooth, zero-entropy cosmological beginning is the expected outcome
        under the compression measure --- not a fine-tuned accident. No
        residual quantum irregularities are required to produce a rich
        structured universe.
\end{enumerate}

\begin{openbox}
\textbf{\textcolor{gold}{Open problems}}

\vspace{0.4em}

\textbf{Fermionic exclusion.} The natural emergence of the Pauli exclusion
principle from description-length explosion upon observer merger is sketched
but not yet demonstrated computationally. This is the current experimental
frontier.

\vspace{0.4em}

\textbf{Dirac unification.} The rigorous proof that $\Cs + \Cg$ is the scale
decomposition of $\mathrm{Tr}(f(D/\Lambda))$ awaits verification. This is
the primary mathematical open problem.

\vspace{0.4em}

\textbf{Particle spectrum.} Identifying the self-consistent structural filter
whose lognormal parameters constitute a derivation of the observed particle
spectrum remains open.
\end{openbox}

% ─────────────────────────────────────────────────────────────────────────────
\section{What this changes}

The \iame{} framework does not propose a new particle or a new symmetry. It
proposes a different answer to the question of what physics is for. Physics
is not a description of a universe that exists independently of observers.
It is the structure of what a finite, compressible observer is expected to
find when it looks. The laws of physics are not external impositions on
reality --- they are the stable invariants of the compressible structures
that observers are made of. Mathematics is not unreasonably effective --- it
is exactly as effective as you would expect once you notice that observers
and mathematics are carved from the same compressible structure.

Time does not flow. It is an ordering relation internal to the observer. The
universe does not expand into anything. Expansion is what entropy increase
looks like when read geometrically. Singularities are not catastrophes. They
are the simplest state. And the universe you experience is not your
container. It is your interior.

\vspace{1em}

\begin{center}
\itshape
The observer is not inside the universe. The universe is inside the observer.
\end{center}

\vspace{2em}

\noindent{\color{rulecolor}\rule{\linewidth}{0.5pt}}

\vspace{0.75em}

{\small\color{inkmuted}
This introduction accompanies four technical papers: Paper~0 (axiomatic
foundations, 2001), Paper~II (black hole singularities, 2010), Paper~III
(spacetime expansion, 2012), and Paper~I (the full \iame{} framework).
Proof-of-concept simulations demonstrating emergent inertia, wave
interference, and gravitational attraction from pure compression optimisation
are provided in the supplementary material.}

\end{document}
